% ThermalQualityRevision20261023-v1
% QualityRevision20261023
\hypersetup{bookmarksdepth=3}
\begingroup
\renewcommand{\Needspace}[1]{\par\begingroup
\dimen0=#1\relax\vskip0pt plus\dimen0\penalty-100
\vskip0pt plus-\dimen0\vskip\dimen0\penalty9999\vskip-\dimen0
\vskip0pt\endgroup}
\tcbset{breakable/.default=false}
\chapter{热力学第二定律}
% Editable TikZ and prose: state entropy versus actual process.
\providecommand{\ThermalFigure}[3]{\begin{qualityfigure}
\begin{tikzpicture}[>=Latex,font=\small,x=.8cm,y=.8cm]#1\end{tikzpicture}
\figcap{#2}\label{#3}\end{qualityfigure}}
\newcommand{\ThQTwoPaths}{\ThermalFigure{
 \draw[thick](0,0)rectangle(4,2);\draw[line width=2pt](2,0)--(2,2);
 \node at(1,1){$T_0,V_0$};\node at(3,1){真空};
 \draw[->,thick](4.5,1)--(5.5,1);\draw[thick](6,0)rectangle(10,2);\node at(8,1){$T_f,2V_0$};
 \node[below,align=center]at(5,-.4){甲：热库 $T_0$，可逆膨胀\quad 乙：隔热，撤去隔板};
}{同一初态到体积加倍的两种实现方式；末温及熵产生由过程分别判断。}{ch16:quality:paths-q}}
\newcommand{\ThATwoPaths}{\ThermalFigure{
 \node[draw,minimum width=2.4cm,minimum height=1cm] (A)at(0,0){$(T_0,V_0)$};
 \node[draw,minimum width=2.4cm,minimum height=1cm] (B)at(7,0){$(T_0,2V_0)$};
 \draw[->,thick,DeepBlue](A.north)..controls(2,2)and(5,2)..(B.north)node[midway,above]{甲：可逆等温，$Q_{\rm rev}$};
 \draw[->,thick](A.south)..controls(2,-2)and(5,-2)..(B.south)node[midway,below]{乙：自由膨胀，$Q=0$};
}{气体的熵变由相同端点决定，总熵产生还须计入热库。}{ch16:quality:paths-a}
图~\ref{ch16:quality:paths-a}的两条箭头连接同一对平衡态，所以气体的 $\Delta S$ 相同。乙过程不能用真实热量 $Q=0$ 代替上方辅助可逆路径的 $Q_{\rm rev}$；两种路径的环境熵变不同，才使总熵产生不同。
}
\newcommand{\ThFiniteWorkBase}[1]{
 \node[draw,minimum width=2.4cm,minimum height=.9cm](H)at(0,2){$C_1,T_h$};
 \node[draw,minimum width=2.4cm,minimum height=.9cm](L)at(0,-2){$C_2,T_c$};
 \node[draw,circle,minimum size=1cm](E)at(0,0){热机};
 \draw[->,thick](H)--node[right]{$\delta Q_h$}(E);\draw[->,thick](E)--node[right]{$\delta Q_c$}(L);
 \draw[->,DeepBlue,thick](E)--(3,0)node[right]{$\delta W$};#1}
\newcommand{\ThQFiniteWork}{\ThermalFigure{\ThFiniteWorkBase{
 \node[right,align=left]at(4,1){两端是有限物体\\温度均随取热、放热改变};
}}{整体隔热，只向外输出功；两物体最终达到共同温度。}{ch16:quality:finite-q}}
\newcommand{\ThAFiniteWork}{\ThermalFigure{
 \draw[->](0,0)--(6,0)node[right]{过程进度};\draw[->](0,0)--(0,3.6)node[above]{$T$};
 \draw[DeepBlue,thick](.3,3.1)..controls(2,2.3)and(4,1.8)..(5.5,1.8);
 \draw[thick](.3,.6)..controls(2,1.3)and(4,1.8)..(5.5,1.8);
 \node[left]at(0,3.1){$T_h$};\node[left]at(0,.6){$T_c$};\node[right]at(5.5,1.8){$T_*$（可逆极限）};
}{热物体降温，冷物体升温；可逆运行给出能达到的最低共同终温。}{ch16:quality:finite-a}
图~\ref{ch16:quality:finite-a}中的共同终温不是预先给定的热库温度。终温越低，留在两物体中的总内能越少，输出功越大；因此先用总熵不减确定终温下界，再由能量守恒求最大功。
}
\newcommand{\ThCoolingBase}[1]{
 \node[draw,minimum width=2.6cm,minimum height=.9cm](C)at(0,0){物体 $C,T$};
 \node[draw,circle,minimum size=1cm](R)at(4,0){制冷机};
 \node[draw,minimum width=2.6cm,minimum height=.9cm](H)at(8,0){环境 $T_0$};
 \draw[->,thick](C)--node[above]{$\delta Q_c$}(R);\draw[->,thick](R)--node[above]{$\delta Q_h$}(H);
 \draw[->,DeepBlue,thick](4,-1.7)--(R)node[midway,right]{$\delta W$};#1}
\newcommand{\ThQCooling}{\ThermalFigure{\ThCoolingBase{}}{制冷机循环复原，有限物体从 $T_0$ 冷却至 $T_1$。}{ch16:quality:cooling-q}}
\newcommand{\ThACooling}{\ThermalFigure{\ThCoolingBase{
 \node[below]at(0,-.9){$\dd T<0,\quad\delta Q_c=-C\dd T>0$};
 \node[above]at(4,1){可逆元循环：$\delta Q_h/T_0=\delta Q_c/T$};
}}{按物体瞬时温度逐次计算取热与元功，再沿降温过程积分。}{ch16:quality:cooling-a}
图~\ref{ch16:quality:cooling-a}中的热量箭头均取正的大小。物体每降低 $\dd T<0$ 放出 $-C\dd T$，环境收到的热量还包含输入功；可逆熵平衡给出 $\delta W=\delta Q_c(T_0/T-1)$，随 $T$ 改变，不能把终态制冷系数用于整个过程。
}
\newcommand{\ThQCarnot}{\ThermalFigure{
 \node[draw](H)at(0,2){热库 $T_1$};\node[draw,circle](E)at(0,0){循环工质};\node[draw](C)at(0,-2){热库 $T_2$};
 \draw[->,thick](H)--node[right]{$Q_1$}(E);\draw[->,thick](E)--node[right]{$Q_2$}(C);
 \draw[->,DeepBlue,thick](E)--(3,0)node[right]{$W$};\node[right]at(4,1.5){$T_1>T_2$};
}{只与两个热库交换热量的循环热机；先约定热量大小与方向。}{ch16:quality:carnot-q}}
\newcommand{\ThQSecondLaw}{\ThermalFigure{
 \node[draw](H)at(0,1.5){单一热库};\node[draw,circle](E)at(0,0){循环装置};
 \draw[->](H)--node[right]{$Q$}(E);\draw[->](E)--(2.3,0)node[right]{$W$};
 \node[below]at(0,-1){开尔文情境};
 \node[draw](C)at(6,0){低温热库 $T_c$};\node[draw](T)at(6,2){高温热库 $T_h$};
 \draw[->,thick](C)--node[right]{$Q$}(T);\node[below]at(6,-1){克劳修斯情境};
}{两种假设过程：完全把单库热量变成功，或不附带其他变化地逆温差传热。}{ch16:quality:second-law-q}}
\newcommand{\ThQIrreversibility}{\ThermalFigure{
 \draw[->](0,0)--(4.7,0)node[right]{$V$};\draw[->](0,0)--(0,3)node[above]{$p$};
 \draw[DeepBlue,thick,domain=1:4,samples=60]plot(\x,{2.6/\x});
 \draw[thick,domain=1:4,samples=60]plot(\x,{2.6/\x^1.5});
 \fill(1,2.6)circle(1.5pt)node[above right]{$A$};
 \node[right]at(4.8,2.2){等温与可逆绝热};
 \node[right]at(4.8,1.2){还需分别讨论有限温差传热、快速膨胀};
}{状态曲线的几何性质与过程能否自然发生属于不同问题，需要分别核查条件。}{ch16:quality:irreversibility-q}}
\newcommand{\ThQThreeBodies}{\ThermalFigure{
 \foreach \x/\lab/\temp in {0/1/300,3/2/300,6/3/100}{
 \draw[thick](\x,0)rectangle(2+\x,1.5);\node at(\x+1,1){物体 $\lab$};\node at(\x+1,.45){$c,\ \temp\,\mathrm K$};}
 \draw[dashed](-.4,-.4)rectangle(8.4,2);\node[above]at(4,2.1){外界不供热、不做功；内部可用循环热机};
}{三个物体组成隔离整体，最高温度由总能量与总熵共同限制。}{ch16:quality:three-q}}
\newcommand{\ThQSurfaceEngine}{\ThermalFigure{
 \draw[thick](0,0)rectangle(4,2);\fill[DeepBlue!12](0,0)rectangle(2.7,2);
 \draw[line width=2pt](2.7,0)--(2.7,2);\draw[->](2.7,2.5)--(3.7,2.5)node[right]{$\dd x$};
 \draw[<->](-.5,0)--node[left]{$\ell$}(-.5,2);\node at(1.3,1){表面 $A$};
 \draw[->,DeepBlue,thick](2.8,1)--(4.7,1)node[right]{维持准静态的外力};
 \node[right,align=left]at(5.4,1.4){界面张力 $\sigma(T)$\\温度由热库控制};
}{可移动边界改变界面面积；外力功与表面张力功的符号应随面积变化判断。}{ch16:quality:surface-q}}



{
\graphicspath{{../../../sample/xes-original/thirtieth-section/}{../../../sample/xes-original/thirtieth-section/assets/}}

% Original vector redraws for Lesson 21.  Text labels deliberately have no
% rectangular background; spacing is achieved by placement, not masking.

\tikzset{
  thermobody/.style={draw=Ink,thick},
  thermoarrow/.style={-{Latex[length=2.2mm]},draw=DeepBlue,very thick},
  thermoguide/.style={draw=TextGray,densely dashed},
  thermolabel/.style={font=\scriptsize,text=Ink,fill=none,inner sep=0pt,outer sep=.8pt}
}


\newcommand{\pistonstage}[4]{%
  \begin{scope}[shift={(#1,0)}]
    \draw[thermobody] (-0.62,-1.20)--(-0.62,1.15) (0.62,-1.20)--(0.62,1.15);
    \draw[thermobody] (-0.62,1.15)--(0.62,1.15);
    \draw[fill=SoftGray,draw=Ink,thick] (-0.54,#2) rectangle (0.54,{#2+0.18});
    \coordinate (Pin) at (0,#2);
    \coordinate (Cen) at (0,-1.78);
    \coordinate (Crank) at ($(Cen)+(#3:0.52)$);
    \draw[Ink,thick] (Pin)--(Crank);
    \draw[Ink,thick] (Cen) circle (0.52);
    \draw[Ink,thick] (Cen)--(Crank);
    \fill[Ink] (Cen) circle (1.1pt) (Crank) circle (1.1pt);
    \node[thermolabel,anchor=north] at (0,-2.42) {#4};
  \end{scope}
}

\newcommand{\ottoCycleFigure}{%
\begin{center}
\begin{tikzpicture}[x=.92cm,y=.92cm]
  \pistonstage{-6.8}{0.05}{220}{进气}
  \pistonstage{-5.05}{0.52}{300}{压缩}
  \pistonstage{-3.30}{0.72}{80}{点火}
  \pistonstage{-1.55}{-0.10}{155}{做功}
  \pistonstage{0.20}{0.18}{235}{排气}
  \draw[thermoarrow] (-7.85,0.85)--(-7.10,0.85);
  \node[thermolabel,above] at (-7.48,0.85) {进气};
  \draw[thermoarrow] (0.53,0.85)--(1.25,0.85);
  \node[thermolabel,above] at (0.88,0.85) {排气};
  \draw[thermoarrow] (-6.8,1.85)--(-6.8,1.25);
  \node[thermolabel,above right] at (-6.8,1.85) {火花塞};
  \draw[DeepBlue,thick] (-3.53,1.22)--(-3.30,1.57)--(-3.10,1.22)
    (-3.45,1.35)--(-3.28,1.80)--(-3.18,1.37);

  \begin{scope}[shift={(2.35,-2.15)},x=.9cm,y=.75cm]
    \draw[axis] (0,0)--(5.4,0) node[right] {$V$};
    \draw[axis] (0,0)--(0,4.9) node[above] {$P$};
    \coordinate (O) at (.55,.85);
    \coordinate (A) at (4.55,.85);
    \coordinate (B) at (.90,2.35);
    \coordinate (C) at (.90,4.15);
    \coordinate (D) at (4.55,1.75);
    \draw[DeepBlue,very thick,postaction={decorate},decoration={markings,
      mark=at position .35 with {\arrow{Latex}},mark=at position .78 with {\arrow{Latex}}}]
      (O)--(A)--(B)--(C) .. controls (1.85,3.55) and (3.20,2.35) .. (D)--(A)--(O);
    \foreach \p/\lab/\pos in {O/O/below left,A/A/below right,B/B/left,C/C/above left,D/D/above right}{
      \fill[Ink] (\p) circle (1.7pt) node[thermolabel,\pos] {$\lab$};
    }
    \draw[thermoguide] (.90,0)--(.90,4.15) (4.55,0)--(4.55,1.75);
    \node[thermolabel,below] at (.90,0) {$V_2$};
    \node[thermolabel,below] at (4.55,0) {$V_1$};
    \draw[thermoarrow] (1.25,2.75)--(1.25,3.65);
    \node[thermolabel,right] at (1.25,3.20) {$Q_h$};
    \draw[thermoarrow] (4.85,1.70)--(5.30,1.35);
    \node[thermolabel,right] at (5.18,1.48) {$Q_c$};
  \end{scope}
\end{tikzpicture}
\end{center}\nopagebreak[2]}

\newcommand{\carnotCycleFigure}{%
\begin{center}
\begin{tikzpicture}[x=1.08cm,y=.95cm]
  \draw[axis] (0,0)--(7.4,0) node[right] {$V$};
  \draw[axis] (0,0)--(0,5.0) node[above] {$P$};
  \draw[TextGray,thick,domain=.8:6.7,samples=100] plot (\x,{4.55/(\x^.43)});
  \draw[TextGray,thick,domain=.8:6.7,samples=100] plot (\x,{2.30/(\x^.43)});
  \coordinate (P1) at (1.35,4.00);
  \coordinate (P2) at (4.35,2.48);
  \coordinate (P3) at (5.55,1.10);
  \coordinate (P4) at (2.05,1.70);
  \draw[DeepBlue,very thick,postaction={decorate},decoration={markings,
    mark=at position .12 with {\arrow{Latex}},mark=at position .38 with {\arrow{Latex}},
    mark=at position .63 with {\arrow{Latex}},mark=at position .88 with {\arrow{Latex}}}]
    (P1) .. controls (2.25,3.38) and (3.35,2.76) .. (P2)
    .. controls (4.75,1.83) and (5.10,1.38) .. (P3)
    .. controls (4.35,1.20) and (3.00,1.45) .. (P4)
    .. controls (1.90,2.62) and (1.62,3.36) .. cycle;
  \foreach \p/\lab/\pos in {P1/1/above left,P2/2/above right,P3/3/below right,P4/4/below left}{
    \fill[Ink] (\p) circle (1.8pt) node[thermolabel,\pos] {$\lab$};
  }
  \node[thermolabel,anchor=west] at (5.55,3.15) {$T_1$ 等温线};
  \node[thermolabel,anchor=west] at (4.85,.55) {$T_2$ 等温线};
  \node[thermolabel] at (3.0,3.18) {$Q_1$};
  \node[thermolabel] at (3.75,1.15) {$Q_2$};
  \node[thermolabel,DeepBlue] at (4.05,4.35) {$T_1>T_2$};
\end{tikzpicture}
\end{center}\nopagebreak[2]}

\newcommand{\heatdevicepanel}[5]{%
  \begin{scope}[shift={(#1,0)}]
    \draw[thermobody,fill=SoftGray] (-1.40,2.15) rectangle (1.40,2.62);
    \draw[thermobody,fill=SoftGray] (-1.40,-2.62) rectangle (1.40,-2.15);
    \node[thermolabel] at (0,2.38) {高温热库 $T_h$};
    \node[thermolabel] at (0,-2.38) {低温热库 $T_c$};
    \draw[thermobody,fill=PaleBlue] (0,0) circle (.64);
    \node[thermolabel] at (0,0) {#2};
    #3 #4 #5
  \end{scope}
}

\newcommand{\heatDevicesFigure}{%
\begin{center}
\begin{tikzpicture}[x=1cm,y=.82cm]
  \heatdevicepanel{-5.7}{热机}{\draw[thermoarrow] (0,2.12)--(0,.70) node[midway,right,thermolabel] {$Q_h$};}{\draw[thermoarrow] (0,-.70)--(0,-2.12) node[midway,right,thermolabel] {$Q_c$};}{\draw[thermoarrow] (.68,0)--(1.55,0) node[midway,above,thermolabel] {$W$};}
  \heatdevicepanel{-1.9}{“完美”热机}{\draw[thermoarrow] (0,2.12)--(0,.70) node[midway,right,thermolabel] {$Q_h$};}{}{\draw[thermoarrow] (.68,0)--(1.55,0);\node[thermolabel] at (1.15,-1.05) {$W=Q_h$};}
  \heatdevicepanel{1.9}{制冷机}{\draw[thermoarrow] (0,.70)--(0,2.12) node[midway,right,thermolabel] {$Q_h$};}{\draw[thermoarrow] (0,-2.12)--(0,-.70) node[midway,right,thermolabel] {$Q_c$};}{\draw[thermoarrow] (1.55,0)--(.68,0) node[midway,above,thermolabel] {$W$};}
  \heatdevicepanel{5.7}{“完美”制冷机}{\draw[thermoarrow] (0,.70)--(0,2.12);}{\draw[thermoarrow] (0,-2.12)--(0,-.70);}{}
  \node[thermolabel] at (-5.7,-3.05) {(a) 真实热机};
  \node[thermolabel] at (-1.9,-3.05) {(b) 第二类永动机};
  \node[thermolabel] at (1.9,-3.05) {(c) 真实制冷机};
  \node[thermolabel] at (5.7,-3.05) {(d) 无功制冷机};
\end{tikzpicture}
\end{center}\nopagebreak[2]}

\newcommand{\portraitOne}[5]{%
  \begin{scope}[shift={(#1,0)}]
    \draw[fill=Ink,draw=Ink,rounded corners=1pt] (-1.62,-2.15) rectangle (1.62,2.25);
    \fill[#2] (-1.18,-2.10) .. controls (-1.02,-1.00) and (-.75,-.63) .. (-.48,-.50)
      --(.48,-.50) .. controls (.82,-.72) and (1.07,-1.15) .. (1.18,-2.10)--cycle;
    \fill[SoftGray] (0,.35) ellipse (.68 and 1.02);
    #3
    \draw[Ink,thick] (-.30,.47)--(-.08,.50) (.10,.50)--(.32,.47);
    \fill[Ink] (-.18,.49) circle (.04) (.18,.49) circle (.04);
    \draw[Ink,thin] (0,.37)--(-.05,.05)--(.06,.00);
    \draw[Ink,thin] (-.24,-.17) .. controls (0,-.28) and (.18,-.22) .. (.28,-.14);
    #4
    \node[font=\small\bfseries,anchor=north] at (0,-2.38) {#5};
  \end{scope}
}

\newcommand{\foundersPortraitFigure}{%
\begin{center}
\begin{tikzpicture}
  \portraitOne{-4.6}{TextGray}{\fill[Ink] (-.70,.82) .. controls (-.48,1.58) and (.55,1.58) .. (.72,.74) .. controls (.28,1.10) and (-.25,1.15) .. cycle;}{\draw[SoftGray,very thick] (-.40,-.62)--(-.15,-.88)--(0,-.70)--(.18,-.88)--(.42,-.62);}{萨迪·卡诺}
  \portraitOne{0}{RuleGray}{\fill[SoftGray] (-.60,.92) arc[start angle=180,end angle=0,x radius=.60,y radius=.50];}{\fill[SoftGray] (-.52,-.18) .. controls (-.62,-.95) and (-.35,-1.28) .. (0,-1.48) .. controls (.36,-1.25) and (.62,-.92) .. (.50,-.18) .. controls (.18,-.42) and (-.18,-.42) .. cycle;}{开尔文勋爵}
  \portraitOne{4.6}{TextGray}{\fill[Ink] (-.66,.88) .. controls (-.38,1.42) and (.48,1.34) .. (.68,.70) .. controls (.30,1.02) and (-.15,1.02) .. cycle;}{\draw[Ink,very thick] (-.34,-.12) .. controls (0,-.02) and (.18,-.04) .. (.36,-.12);}{鲁道夫·克劳修斯}
\end{tikzpicture}
\end{center}\nopagebreak[2]}

\newcommand{\entropyPathFigure}{%
\begin{center}
\begin{tikzpicture}[x=1cm,y=.9cm]
  \begin{scope}[shift={(-3.8,-1.4)}]
    \draw[axis] (0,0)--(5.0,0) node[right] {$V$};
    \draw[axis] (0,0)--(0,4.3) node[above] {$P$};
    \coordinate (A) at (1.0,1.0);
    \coordinate (B) at (1.0,3.55);
    \coordinate (C) at (4.25,.85);
    \draw[DeepBlue,very thick,->] (A)--(B);
    \draw[DeepBlue,very thick,->] (B) .. controls (2.0,2.95) and (3.0,1.70) .. (C);
    \fill[Ink] (A) circle (1.7pt) node[thermolabel,left] {$(V_0,T_0)$};
    \fill[Ink] (B) circle (1.7pt) node[thermolabel,above left] {$(V_0,T)$};
    \fill[Ink] (C) circle (1.7pt) node[thermolabel,right] {$(V,T)$};
    \node[thermolabel,anchor=west] at (1.15,2.25) {等容升温};
    \node[thermolabel,rotate=-28] at (2.90,2.05) {等温膨胀};
  \end{scope}
  \node[draw=DeepBlue,rounded corners=2pt,thick,align=left,inner sep=7pt] at (4.0,.2)
    {$\displaystyle \Delta S=n c_v\ln\frac{T}{T_0}+nR\ln\frac{V}{V_0}$\\[5pt]
     熵是状态量：可任选一条可逆路径计算。};
\end{tikzpicture}
\end{center}\nopagebreak[2]}

\newcommand{\boltzmannMemorialFigure}{%
\begin{center}
\begin{tikzpicture}[x=1cm,y=1cm]
  \fill[TextGray!28] (-4.6,-2.5) rectangle (4.6,1.65);
  \fill[TextGray!28] (-4.6,1.60) arc[start angle=180,end angle=0,x radius=4.6,y radius=2.0]--cycle;
  \draw[Ink,thick] (-4.6,-2.5)--(-4.6,1.6) arc[start angle=180,end angle=0,x radius=4.6,y radius=2.0]--(4.6,-2.5)--cycle;
  \node[font=\large\bfseries] at (0,2.25) {$S=k_B\ln\Omega$};
  \fill[Ink] (-1.35,-2.48) .. controls (-1.05,-1.35) and (-.78,-.78) .. (-.48,-.55)
    --(.48,-.55) .. controls (.82,-.86) and (1.12,-1.45) .. (1.35,-2.48)--cycle;
  \fill[SoftGray] (0,.10) ellipse (.80 and 1.12);
  \fill[Ink] (-.75,.54) .. controls (-.42,1.36) and (.52,1.36) .. (.78,.48)
    .. controls (.36,.84) and (-.20,.88) .. cycle;
  \draw[Ink,thick] (-.40,.25)--(-.12,.28) (.12,.28)--(.40,.25);
  \fill[Ink] (-.22,.26) circle (.045) (.22,.26) circle (.045);
  \draw[Ink,thin] (0,.15)--(-.05,-.18)--(.06,-.15);
  \fill[SoftGray] (-.55,-.20) .. controls (-.62,-1.02) and (-.30,-1.35) .. (0,-1.52)
    .. controls (.33,-1.35) and (.64,-1.00) .. (.55,-.20) .. controls (.16,-.45) and (-.16,-.45) .. cycle;
  \draw[thermoguide] (-3.8,-1.95)--(-2.6,-.25)--(-3.3,.70);
  \draw[thermoguide] (3.8,-1.95)--(2.6,-.25)--(3.3,.70);
\end{tikzpicture}
\end{center}\nopagebreak[2]}

\newcommand{\carnotProofDiagram}{%
\begin{center}
\begin{tikzpicture}[x=1cm,y=1cm]
  \node[draw=Ink,fill=SoftGray,minimum width=3.0cm,minimum height=.65cm] (H) at (0,2.0) {高温热库 $T_h$};
  \node[draw=DeepBlue,fill=PaleBlue,minimum width=2.2cm,minimum height=.8cm] (E) at (0,0) {可逆热机};
  \node[draw=Ink,fill=SoftGray,minimum width=3.0cm,minimum height=.65cm] (C) at (0,-2.0) {低温热库 $T_c$};
  \draw[thermoarrow] (H)--node[right,thermolabel] {$Q_h$}(E);
  \draw[thermoarrow] (E)--node[right,thermolabel] {$Q_c$}(C);
  \draw[thermoarrow] (E)--(2.25,0) node[right,thermolabel] {$W=Q_h-Q_c$};
  \node[draw=DeepBlue,rounded corners=2pt,align=left,inner sep=6pt] at (5.4,0)
    {$\Delta S_{\rm cyc}=0$\\[2pt]
     $\displaystyle \frac{Q_h}{T_h}-\frac{Q_c}{T_c}=0$\\[4pt]
     $\displaystyle \eta=1-\frac{T_c}{T_h}$};
\end{tikzpicture}
\end{center}\nopagebreak[2]}

\newcommand{\secondLawEntropyDiagram}{%
\begin{center}
\begin{tikzpicture}[x=1cm,y=1cm]
  \node[draw=RuleGray,rounded corners=2pt,minimum width=5.4cm,minimum height=2.25cm,align=center] at (-3.5,0)
    {单一热库取热并全部变成功\\[4pt]$\Delta S_{\rm tot}=-Q/T<0$\\[2pt]\color{DeepBlue}不可能};
  \node[draw=RuleGray,rounded corners=2pt,minimum width=5.4cm,minimum height=2.25cm,align=center] at (3.5,0)
    {无补偿地把热从 $T_c$ 送往 $T_h$\\[4pt]
     $\displaystyle\Delta S_{\rm tot}=\frac{Q}{T_h}-\frac{Q}{T_c}<0$\\[2pt]\color{DeepBlue}不可能};
\end{tikzpicture}
\end{center}\nopagebreak[2]}

\newcommand{\irreversibleProcessesDiagram}{%
\begin{center}
\begin{tikzpicture}[x=1cm,y=1cm]
  \begin{scope}[shift={(-4.3,0)}]
    \draw[thermobody] (-2,-1.2) rectangle (2,1.2);
    \draw[thermobody] (0,-1.2)--(0,1.2);
    \fill[PaleBlue] (-2,-1.2) rectangle (0,1.2);
    \node at (-1,0) {气体};
    \draw[thermoarrow] (-.25,0)--(1.35,0);
    \node[thermolabel] at (0,-1.55) {快速自由膨胀：$\Delta S=nR\ln(V_f/V_i)>0$};
  \end{scope}
  \begin{scope}[shift={(4.3,0)}]
    \draw[thermobody,fill=SoftGray] (-2,-1.1) rectangle (-.35,1.1);
    \draw[thermobody,fill=PaleBlue] (.35,-1.1) rectangle (2,1.1);
    \node at (-1.18,0) {$T_h$};
    \node at (1.18,0) {$T_c$};
    \draw[thermoarrow] (-.2,0)--(.2,0);
    \node[thermolabel] at (0,-1.55) {$\Delta S=Q(1/T_c-1/T_h)>0$};
  \end{scope}
\end{tikzpicture}
\end{center}\nopagebreak[2]}

\newcommand{\threeBodiesDiagram}{%
\begin{center}
\begin{tikzpicture}[x=1cm,y=1cm]
  \foreach \x/\t/\lab in {-4/300/{物体1},0/300/{物体2},4/100/{物体3}}{
    \draw[thermobody,fill=SoftGray] (\x-1.05,-.55) rectangle (\x+1.05,.55);
    \node at (\x,.15) {\lab};
    \node[thermolabel] at (\x,-.25) {$\t\,\mathrm K$};
  }
  \draw[thermoarrow,<->] (-2.85,0)--(-1.15,0);
  \draw[thermoarrow,<->] (1.15,0)--(2.85,0);
  \node[thermolabel,above] at (-2.0,0) {可逆热机};
  \node[thermolabel,above] at (2.0,0) {可逆热机};
  \node[draw=DeepBlue,rounded corners=2pt,align=center,inner sep=6pt] at (0,-2.05)
    {$T_{\max}+2T=700\,\mathrm K$\\[2pt]$T_{\max}T^2=300^2\!\times100\,\mathrm K^3$};
\end{tikzpicture}
\end{center}\nopagebreak[2]}

\newcommand{\surfaceEngineDiagram}{%
\begin{center}
\begin{tikzpicture}[x=1cm,y=1cm]
  \begin{scope}[shift={(-3.4,0)}]
    \draw[axis] (0,0)--(5.0,0) node[right] {$A$};
    \draw[axis] (0,0)--(0,4.0) node[above] {$\sigma$};
    \coordinate (A) at (1.0,1.45);
    \coordinate (B) at (4.1,1.45);
    \coordinate (C) at (4.1,3.05);
    \coordinate (D) at (1.0,3.05);
    \draw[DeepBlue,very thick,postaction={decorate},decoration={markings,
      mark=at position .12 with {\arrow{Latex}},mark=at position .37 with {\arrow{Latex}},
      mark=at position .62 with {\arrow{Latex}},mark=at position .87 with {\arrow{Latex}}}]
      (A)--(B)--(C)--(D)--cycle;
    \node[thermolabel,left] at (0,1.45) {$\sigma_h$};
    \node[thermolabel,left] at (0,3.05) {$\sigma_c$};
    \node[thermolabel,below] at (1.0,0) {$A_1$};
    \node[thermolabel,below] at (4.1,0) {$A_2$};
    \node[thermolabel] at (2.55,2.22) {$W_{\rm out}$};
  \end{scope}
  \node[draw=DeepBlue,rounded corners=2pt,align=left,inner sep=6pt] at (4.4,.3)
    {$S=-A\,\mathrm d\sigma/\mathrm dT$\\[3pt]
     $Q_h=T_h\Delta S$\\[3pt]
     $Q_c=T_c\Delta S$\\[3pt]
     $\eta=1-T_c/T_h$};
\end{tikzpicture}
\end{center}\nopagebreak[2]}

% Shared coordinates for mechanical contact and thermodynamic paths.
\tikzset{revcycle/.style={DeepBlue,thick,postaction={decorate},decoration={markings,mark=at position .55 with {\arrow{Latex}}}}}
\newcommand{\pistonStageRevision}[4]{\begin{scope}[shift={(#1,0)}]
\pgfmathsetmacro{\pinY}{-2+.4*cos(#2)+sqrt(1.5*1.5-.4*.4*sin(#2)*sin(#2))}
\coordinate(pin)at(0,\pinY);\coordinate(crank)at({.4*sin(#2)},{-2+.4*cos(#2)});
\fill[DeepBlue!7](-.62,\pinY+.12)rectangle(.62,.75);
\draw[thick](-.62,-1.2)--(-.62,.75) (.62,-1.2)--(.62,.75) (-.25,.75)--(.25,.75);
\ifnum#4=1 \draw[thick](-.62,.55)--(-.25,.4);\else\draw[thick](-.62,.75)--(-.25,.75);\fi
\ifnum#4=5 \draw[thick](.25,.4)--(.62,.55);\else\draw[thick](.25,.75)--(.62,.75);\fi
\draw[thick,fill=SoftGray](-.62,\pinY-.06)rectangle(.62,\pinY+.12);
\draw[thick](pin)--(crank)--(0,-2);\draw[thick](0,-2)circle(.4);\fill(pin)circle(1.2pt)(crank)circle(1.2pt)(0,-2)circle(1.2pt);
\draw[thick](0,1.03)--(0,.59);\draw[thick](-.09,.64)--(.09,.64);
\ifnum#4=3 \draw[orange!85!black,thick](-.14,.52)--(.05,.38)--(-.06,.24)--(.13,.12);\fi
\ifnum#4=1 \draw[->,DeepBlue,thick](-1.05,.95)--(-.45,.95)--(-.45,.62);\fi
\ifnum#4=5 \draw[->,DeepBlue,thick](.45,.62)--(.45,.95)--(1.05,.95);\fi
\ifnum#4=3\else
\ifnum#4=2 \draw[->](.87,-.8)--(.87,-.1);\else
\ifnum#4=5 \draw[->](.87,-.8)--(.87,-.1);\else\draw[->](.87,-.1)--(.87,-.8);\fi\fi\fi
\node[font=\small]at(0,-2.8){#3};\end{scope}}
\renewcommand{\ottoCycleFigure}{\begin{center}\begin{tikzpicture}[font=\small]
\pistonStageRevision{0}{90}{进气（口开）}{1}
\pistonStageRevision{2.65}{270}{压缩}{2}
\pistonStageRevision{5.3}{0}{点火}{3}
\pistonStageRevision{7.95}{90}{做功}{4}
\pistonStageRevision{10.6}{270}{排气（口开）}{5}
\node[anchor=west]at(6.5,1.45){火花塞};\draw[TextGray](6.35,1.42)--(5.3,1.03);
\node[font=\scriptsize]at(5.3,-3.35){相同缸径、曲柄半径与连杆长度；同一位置可属于不同冲程}
;\end{tikzpicture}
\par\medskip\begin{tikzpicture}[x=1.6cm,y=.9cm,font=\small]
\pgfmathsetmacro{\pa}{1.8*pow(1.2/3,1.4)}\pgfmathsetmacro{\pd}{4*pow(1.2/3,1.4)}
\draw[axis](0,0)--(4,0)node[right]{$V$};\draw[axis](0,0)--(0,4.8)node[above]{$P$};
\draw[revcycle,domain=3:1.2,samples=90]plot(\x,{1.8*pow(1.2/\x,1.4)});
\draw[revcycle](1.2,1.8)--(1.2,4);
\draw[revcycle,domain=1.2:3,samples=90]plot(\x,{4*pow(1.2/\x,1.4)});
\draw[revcycle](3,\pd)--(3,\pa);
\draw[guide](1.2,0)--(1.2,1.65) (3,0)--(3,\pa);
\draw[TextGray,dashed,->](1.2,\pa)--(2.5,\pa);\draw[TextGray,->](2.4,\pa-.12)--(1.6,\pa-.12);
\foreach \x/\y/\n/\a in {3/\pa/A/below right,1.2/1.8/B/left,1.2/4/C/above,3/\pd/D/above right,1.2/\pa/O/left}{\fill(\x,\y)circle(1.2pt)node[\a]{$\n$};}
\node[below]at(1.2,0){$V_2$};\node[below]at(3,0){$V_1$};
\node[anchor=west,align=left]at(4.6,3.2){蓝线：封闭空气标准循环\\[6pt]$A\to B$、$C\to D$：绝热\\[6pt]$B\to C$：等容吸热 $Q_h$\\[6pt]$D\to A$：等容放热 $Q_c$};
\node[anchor=west,align=left,font=\scriptsize]at(4.6,.65){灰色往返线仅提示真实进、排气\\不是蓝色闭式循环的一部分};
\end{tikzpicture}\end{center}\nopagebreak[2]}

\renewcommand{\carnotCycleFigure}{\begin{center}\begin{tikzpicture}[x=1.45cm,y=1cm,font=\small]
\pgfmathsetmacro{\cold}{4/pow(1.5,2/3)}
\draw[axis](0,0)--(4.4,0)node[right]{$V$};\draw[axis](0,0)--(0,4.9)node[above]{$P$};
\draw[TextGray,domain=.9:4.15,samples=100]plot(\x,{4/\x});
\draw[TextGray,domain=.9:4.15,samples=100]plot(\x,{\cold/\x});
\draw[revcycle,domain=1:2.4,samples=85]plot(\x,{4/\x});
\draw[revcycle,domain=2.4:3.6,samples=85]plot(\x,{4*pow(2.4,2/3)/pow(\x,5/3)});
\draw[revcycle,domain=3.6:1.5,samples=85]plot(\x,{\cold/\x});
\draw[revcycle,domain=1.5:1,samples=85]plot(\x,{4/pow(\x,5/3)});
\foreach \x/\y/\n/\a in {1/4/1/above left,2.4/{4/2.4}/2/above right,3.6/{\cold/3.6}/3/below,1.5/{\cold/1.5}/4/below left}{\fill(\x,\y)circle(1.4pt)node[\a]{$\n$};}
\node[anchor=west,align=left]at(5,3.4){高温等温：$PV=nRT_1$\\[5pt]$1\to2$ 吸热 $Q_1$};
\node[anchor=west,align=left]at(5,1.6){低温等温：$PV=nRT_2$\\[5pt]$3\to4$ 放热 $Q_2$};
\draw[TextGray](4.1,{4/4.1})--(4.55,3.4)--(4.8,3.4);
\draw[TextGray](4.1,{\cold/4.1})--(4.65,1.6)--(4.8,1.6);
\end{tikzpicture}\end{center}\nopagebreak[2]}

\renewcommand{\heatdevicepanel}[5]{\begin{scope}[shift={(#1,0)}]
\draw[thermobody,fill=SoftGray](-1.4,2.15)rectangle(1.4,2.62);\draw[thermobody,fill=SoftGray](-1.4,-2.62)rectangle(1.4,-2.15);
\node[thermolabel]at(0,2.38){高温热库 $T_h$};\node[thermolabel]at(0,-2.38){低温热库 $T_c$};
\draw[thermobody](-.66,-.65)rectangle(.66,.65);\node[thermolabel,align=center,text width=1.1cm]at(0,0){#2};#3 #4 #5\end{scope}}
\renewcommand{\foundersPortraitFigure}{\begin{center}\includegraphics[width=.96\linewidth]{thermodynamics-founders.png}\par\smallskip
\makebox[.31\linewidth]{萨迪·卡诺}\makebox[.31\linewidth]{开尔文}\makebox[.31\linewidth]{鲁道夫·克劳修斯}\end{center}\nopagebreak[2]}
\renewcommand{\boltzmannMemorialFigure}{\begin{center}\begin{tikzpicture}
\node[inner sep=0]at(0,0){\includegraphics[width=6.2cm]{boltzmann-memorial.png}};
\node[font=\large]at(0,2.32){$S=k_B\ln\Omega$};\end{tikzpicture}\end{center}\nopagebreak[2]}

\renewcommand{\entropyPathFigure}{\begin{center}\begin{tikzpicture}[font=\small]
\draw[axis](0,0)--(5,0)node[right]{$V$};\draw[axis](0,0)--(0,4.8)node[above]{$P$};
\draw[revcycle](1,1.6)--(1,4);\draw[revcycle,domain=1:4,samples=90]plot(\x,{4/\x});
\foreach \x/\y/\n/\a in {1/1.6/A/left,1/4/B/above,4/1/C/right}{\fill(\x,\y)circle(1.5pt)node[\a]{$\n$};}
\node[anchor=west,align=left]at(6,3.45){$A:(V_0,T_0)$\\[5pt]$B:(V_0,T)$\\[5pt]$C:(V,T)$};
\node[anchor=west,align=left]at(6,1.25){$A\to B$：等容\\[5pt]$B\to C$：等温，$PV=nRT$};
\end{tikzpicture}\par\smallskip
$\displaystyle\Delta S=n c_v\ln\frac{T}{T_0}+nR\ln\frac{V}{V_0}$
\end{center}\nopagebreak[2]}

\renewcommand{\carnotProofDiagram}{\begin{center}\begin{tikzpicture}[font=\small]
\node[draw,minimum width=3cm,minimum height=.6cm](H)at(0,2){高温热库 $T_1$};
\node[draw,minimum width=2.2cm,minimum height=.8cm](E)at(0,0){可逆热机};
\node[draw,minimum width=3cm,minimum height=.6cm](C)at(0,-2){低温热库 $T_2$};
\draw[thermoarrow](H)--node[right]{$Q_1$}(E);\draw[thermoarrow](E)--node[right]{$Q_2$}(C);
\draw[thermoarrow](E.east)--(3.1,0);\node[above]at(2.2,.15){$W$};
\node[anchor=west,align=left]at(4,0){$W=Q_1-Q_2$\\[8pt]$Q_1/T_1=Q_2/T_2$\\[8pt]$\eta=1-T_2/T_1$};
\end{tikzpicture}\end{center}\nopagebreak[2]}

\renewcommand{\irreversibleProcessesDiagram}{\begin{center}\begin{tikzpicture}[font=\small]
\foreach \sx/\phase in {0/0,5/1}{\begin{scope}[shift={(\sx,0)}]
\ifnum\phase=0 \fill[PaleBlue](-1.6,-1)rectangle(0,1);\draw[thick](0,-1)--(0,1);\node at(-.8,0){气体};\node at(.8,0){真空};
\else\fill[PaleBlue](-1.6,-1)rectangle(1.6,1);\draw[guide](0,-1)--(0,1);\node at(.8,0){气体};\fi
\draw[thick](-1.6,-1)rectangle(1.6,1);\end{scope}}
\draw[->,DeepBlue,thick](1.85,0)--(3.15,0)node[midway,above]{撤去隔板};
\node at(2.5,-1.55){绝热自由膨胀：$Q=W=0$，$\Delta S=nR\ln(V_f/V_i)>0$};
\draw[thick](9,-1)rectangle(10.5,1);\draw[thick](11.3,-1)rectangle(12.8,1);
\node at(9.75,0){$T_h$};\node at(12.05,0){$T_c$};\draw[thermoarrow](10.6,0)--(11.2,0);
\node[align=center]at(10.9,-1.6){两恒温热库间传热\\$\Delta S=Q(1/T_c-1/T_h)>0$};
\end{tikzpicture}\end{center}\nopagebreak[2]}

\renewcommand{\threeBodiesDiagram}{\begin{center}\begin{tikzpicture}[font=\small]
\node[draw,minimum width=2.7cm,minimum height=.7cm](Y)at(0,2.2){供热体 $Y$};
\node[draw,minimum width=2.7cm,minimum height=.7cm](X)at(6,2.2){目标体 $X$};
\node[draw,minimum width=2.7cm,minimum height=.7cm](E)at(0,0){可逆热机};
\node[draw,minimum width=2.7cm,minimum height=.7cm](P)at(6,0){可逆热泵};
\node[draw,minimum width=3cm,minimum height=.7cm](Z)at(3,-2.4){低温体 $Z$};
\draw[thermoarrow](Y)--node[left]{$-c\dd Y$}(E);\draw[thermoarrow](P)--node[right]{$c\dd X$}(X);
\draw[thermoarrow](E.east)--node[above]{$\delta W$}(P.west);
\draw[thermoarrow](E.south)--(0,-2.4)--(Z.west);
\draw[thermoarrow](Z.east)--(6,-2.4)--(P.south);
\node[left,align=right]at(-.1,-1.3){热机向 $Z$ 放热};\node[right,align=left]at(6.1,-1.3){热泵从 $Z$ 吸热};
\node at(3,-3.4){初态 $(X,Y,Z)=(300,300,100)\,\mathrm K$；极限末态 $(400,150,150)\,\mathrm K$};
\end{tikzpicture}\end{center}\nopagebreak[2]}

\renewcommand{\surfaceEngineDiagram}{\begin{center}\begin{tikzpicture}[x=1.5cm,y=.9cm,font=\small]
\draw[axis](0,0)--(3.35,0)node[right]{$A$};\draw[axis](0,0)--(0,4.9)node[above]{$\sigma$};
\fill[PaleBlue](.8,1)--(1.4,1)--plot[domain=1.4:2.8,samples=70](\x,{5-pow(2.8/\x,2)})--(1.6,4)--plot[domain=1.6:.8,samples=70](\x,{5-pow(1.6/\x,2)})--cycle;
\draw[revcycle](.8,1)--(1.4,1);\draw[revcycle,domain=1.4:2.8,samples=70]plot(\x,{5-pow(2.8/\x,2)});
\draw[revcycle](2.8,4)--(1.6,4);\draw[revcycle,domain=1.6:.8,samples=70]plot(\x,{5-pow(1.6/\x,2)});
\foreach \x/\y/\n/\a in {.8/1/1/left,1.4/1/2/below right,2.8/4/3/right,1.6/4/4/above left}{\fill(\x,\y)circle(1.4pt)node[\a]{$\n$};}
\node[left]at(0,1){$\sigma_h$};\node[left]at(0,4){$\sigma_c$};
\node[anchor=west,align=left]at(4.2,3.4){$\sigma(T)=\sigma_*-aT^2$\\[6pt]等熵线：$AT=\text{常量}$\\[6pt]$1\to2$：高温扩面\\[6pt]$3\to4$：低温缩面};
\node[anchor=west]at(4.2,1){$W_{\rm out}=-\oint\sigma\dd A>0$};
\end{tikzpicture}\end{center}\nopagebreak[2]}


从某种意义上讲，热力学的发展与人们制造和改进热机的努力是分不开的。最初，人们设想制造一种
不需要能量输入却能够持续工作的热机，但热力学第一定律告诉我们，这样的努力没有任何希望。退而求
其次，人们又希望得到效率为百分之百的热机，即把吸收的热量全部转化为机械功。从能量守恒的角度看，
这种设想并不直接违背热力学第一定律：能量输入与输出可以相等，但这还没有回答过程是否能够发生。
第二定律额外指出，循环装置不能以“只从单一热库取热并把它全部变成功”为唯一效果。
常见气体热机确有膨胀做功、压缩复原和向外放热的过程，但这只能帮助理解具体模型，不能证明一般定律。
例如液体表面缩小时也能对外做功，后面将研究这种不同的工质。

人们在追求尽可能高效的热机时，逐渐发现这一事实背后隐藏着一条普遍的热力学原理，这就是今天的
热力学第二定律。热机、可逆性以及熵，将是本讲讨论的主题。

\section{热力学循环的效率}

热机在工作时循环运行：工质从某个状态出发，经历一系列过程后再次回到初始状态，完成一个热力学循环。
对在两个热库之间工作的普通热机，工质从高温热库吸热，还向低温热库放热。把整个循环
从外界吸收的总热量记为 $Q^{\rm in}$，把向外排出的总热量记为 $Q^{\rm out}$。由于循环结束后工质回到
初态，内能改变量为零。这里 $Q^{\rm in},Q^{\rm out}$ 分别是累计吸热、累计放热的正量，
也记为 $Q_h,Q_c$；净输入热量为二者之差，不沿用上一章单段过程中带符号的 $Q^{\rm in}=Q$。
热力学第一定律给出循环的净功
\begin{equation}
  \boxed{W^{\rm out}=Q^{\rm in}-Q^{\rm out}}.
\end{equation}
若工质仅以体积功与外界交换机械能，且过程可逆，则净功为 $\oint P\dd V$，等于 $P$--$V$ 图上
围成的有向面积，顺时针为正。对于电功、表面功等形式，须使用相应的功表达式，不能一概用 $P$--$V$ 面积。
只要排出的热量不为零，热机就没有把吸收的全部热量变成功。定义热机
效率为对外净功与吸热量的比值：
\begin{equation}
  \boxed{\eta=\frac{W^{\rm out}}{Q^{\rm in}}
  =1-\frac{Q^{\rm out}}{Q^{\rm in}}}.
\end{equation}

最常见的内燃机可用奥托热机近似描述。图~\ref{ch16:fig:s30-21-1}把活塞的工作过程与相应的 $P$--$V$ 循环放在一起。
从活塞运动看，一个循环可分为进气、压缩、点火、做功和排气几个环节。首先关闭排气口、打开进气口，
空气和汽油的混合气体被吸入气缸；自然吸气发动机的进气压强接近一个大气压，涡轮增压发动机则更高。
活塞到达最下端后进气口关闭，活塞因惯性向上运动并压缩混合气体。压缩过程通常很快，虽然气体温度
升高，但与气缸的热量交换可以忽略，因而可近似为绝热压缩。气体被完全压缩后，火花塞点燃高压油气
混合物，燃烧很快，活塞来不及移动，气体体积近似不变而温度、压强急剧升高。随后高温高压气体推动
活塞向下运动，进入做功冲程；这一快速膨胀过程也可近似为绝热膨胀。活塞到达最下端后排气口打开，
废气被排出，装置完成一个循环。

\par\noindent\begin{minipage}{\linewidth}
\ottoCycleFigure
\diagramcaption{上方是四冲程发动机的五个动作阶段，下方是空气标准奥托循环。点火单列显示，不是额外增加一个活塞冲程。}\label{ch16:fig:s30-21-1}
\end{minipage}\par\medskip

最小、最大气缸体积分别为 $V_2,V_1$。先看真实进排气：$O\to A$ 提示吸气，压强近似不变，
但缸内气体量不断增加，并非定量气体等压膨胀；$A\to O$ 提示排气，气体不断离开气缸。
在真实做功冲程后打开排气门时，废气也会先因内外压差排出。这些开放过程伴随物质携能，
不能套用一份封闭气体的 $nc_v\Delta T$ 来计算排出能量。

为了求一个便于比较的理论效率，现在换用\textbf{空气标准模型}：用同一份定量、定组成的理想气体
代替不断更新的混合气；用可逆等容吸热代替燃烧，用可逆等容放热代替排气散热的效果。
蓝色闭式循环只有四段：$A\to B$ 绝热压缩，$B\to C$ 等容吸热，$C\to D$ 绝热膨胀，
$D\to A$ 等容放热。灰色 $O\leftrightarrow A$ 不是这份封闭工质实际走过的两段。

在这一理想化循环中，只有 $B\to C$ 的等容过程吸热。设气缸内气体可视为总物质的量为 $n$、定容
摩尔热容 $c_v$ 在整个温区内近似恒定的理想气体，并忽略摩擦、漏热和燃烧成分变化，则
\[
  Q_h=n c_v(T_C-T_B).
\]
类似地，只有 $D\to A$ 的等容过程向外排热，排热量为
\[
  Q_c=n c_v(T_D-T_A).
\]
所以
\[
  \eta=1-\frac{Q_c}{Q_h}
  =1-\frac{T_D-T_A}{T_C-T_B}.
\]
再利用两段绝热过程
\[
  T_A V_1^{\gamma-1}=T_B V_2^{\gamma-1},
  \qquad
  T_D V_1^{\gamma-1}=T_C V_2^{\gamma-1},
\]
令压缩比 $r=V_1/V_2>1$，两式分别给出 $T_B=r^{\gamma-1}T_A$、
$T_C=r^{\gamma-1}T_D$。相减而不是直接约去温度，可得
\[
 T_C-T_B=r^{\gamma-1}(T_D-T_A),\qquad
 \frac{Q_c}{Q_h}=\frac{T_D-T_A}{T_C-T_B}=r^{1-\gamma}.
\]
于是得到奥托循环效率与发动机几何参数的关系
\begin{equation}
  \boxed{\eta=1-\left(\frac{V_2}{V_1}\right)^{\gamma-1}}.
\end{equation}
这说明在理想模型中，压缩比 $V_1/V_2$ 越大，循环效率越高；真实发动机还会受到爆震、传热、摩擦和
材料强度等条件限制。

\section{卡诺循环}

在所有热力学循环中，对理论发展最重要的是卡诺设计的卡诺热机。它来自人们寻找最高可能效率
热机的努力；当理论上确定这种热机的效率以后，人们才认识到其背后隐藏着一条基本的热力学定律。
以下示范理想气体循环的计算时，仍取定量气体、恒定 $c_v,c_p$；卡诺定理本身并不限于这一模型。

\par\noindent\begin{minipage}{\linewidth}
\carnotCycleFigure
\diagramcaption{卡诺循环的 $P$--$V$ 图。$1\to2$、$3\to4$ 为等温过程，$2\to3$、$4\to1$ 为绝热过程。}\label{ch16:fig:s30-21-2}
\end{minipage}\par\medskip

如图~\ref{ch16:fig:s30-21-2}所示，卡诺循环由两段绝热过程和两段等温过程组成，两条等温线的温度分别为 $T_1>T_2$。
正是两个热库之间存在温差，才有可能把一部分热能转化为机械功。循环从状态1出发：工质与温度为 $T_1$
的高温热库接触，持续吸热并作等温膨胀，工质对外做功；到达状态2后与热库脱离，继续绝热膨胀，直到
温度降到 $T_2$，来到状态3；然后工质与低温热库接触，受到外界做功而作等温压缩，同时向低温热库
排出热量 $Q_2$；到达状态4后再与热库脱离，作绝热压缩，最后恰好回到初态1。

卡诺可以证明，无论采用何种工作物质，可逆卡诺循环都具有相同效率。先以理想气体为例。设热机
内有 $n$ 摩尔理想气体，高温等温膨胀吸入的热量与低温等温压缩排出的热量分别为
\[
  Q_1=nRT_1\ln\frac{V_2}{V_1},
  \qquad
  Q_2=nRT_2\ln\frac{V_3}{V_4}.
\]
两段绝热过程满足
\[
  T_1V_2^{\gamma-1}=T_2V_3^{\gamma-1},
  \qquad
  T_1V_1^{\gamma-1}=T_2V_4^{\gamma-1}.
\]
两式相除得到 $V_2/V_1=V_3/V_4$。于是效率表达式中所有体积因子恰好抵消，最终得到
\begin{equation}
  \boxed{\eta_{\rm Carnot}=1-\frac{T_2}{T_1}}.
\end{equation}
令人惊讶的是，卡诺热机的效率只与两个热库的热力学温度有关，与热机的几何结构无关；这与奥托
热机效率依赖压缩比形成鲜明对比。

\section{热力学第二定律的早期陈述形式}

卡诺关于热机效率的研究使人们逐渐意识到，有一条基本规律阻止人们制造效率为 $100\%$ 的热机。
这种机器并不直接违反热力学第一定律，却从未能够实现。图~\ref{ch16:fig:s30-21-3}(a)表示真实热机：它从高温热库吸收
$Q_h$，对外做功 $W=Q_h-Q_c$，并向低温热库排出 $Q_c$。图~\ref{ch16:fig:s30-21-3}(b)所示的“完美”热机则只从单一
热库取热并把它全部变成功。开尔文勋爵彻底断绝了这种幻想，他指出：

\begin{quote}\small
不存在这样的装置：当它完成一个循环时，能从一个物体取出热量并把它转换为功，而不发生其他变化
（至少要在其他物体中发生变化）。
\end{quote}

历史上这种机器被称为第二类永动机，因此热力学第二定律的开尔文表述也可理解为：第二类永动机不可能
制成。

\par\noindent\begin{minipage}{\linewidth}
\heatDevicesFigure
\diagramcaption{热机和制冷机能量流。由左至右为普通热机、第二类永动机、普通制冷机和无功制冷机；两种“完美”设想不允许实现。}\label{ch16:fig:s30-21-3}
\end{minipage}\par\medskip

生活经验还指出，热量会自发地从高温物体传向低温物体。冰箱、空调试图把热量从低温部分搬运到高温
环境，必须依赖外部能量供应。在力学中，物体可以凭惯性短时间沿受力的反方向运动，例如竖直上抛物体
仍继续上升，或弹簧振子在回复力作用下仍远离平衡位置；不能据此类比出一种可持续的无补偿逆向传热。
在这里讨论的通常宏观传热中，若不付出其他代价，热量只能自发地从高温传向低温。克劳修斯把这一事实表述为：

\begin{quote}\small
不存在这样的装置：当它完成一个循环时，能使热量从较低温度的热库传到较高温度的热库，而没有补偿
变化。
\end{quote}

从高温向低温的传热过程可以自然发生，而它的逆过程不能独立发生，因此克劳修斯表述也意味着：有限
温差下的热传导是一种不可逆过程。开尔文与克劳修斯的经典英文表述分别为：
\begin{quote}\footnotesize
It is impossible, by means of inanimate material agency, to derive mechanical effect from any portion of
matter by cooling it below the temperature of the coldest of the surrounding objects.

Heat can never pass from a colder to a warmer body without some other change, connected therewith,
occurring at the same time.
\end{quote}

\par\noindent\begin{minipage}{\linewidth}
\foundersPortraitFigure
\diagramcaption{三位热力学奠基人：卡诺研究热机的理想效率，开尔文和克劳修斯提出第二定律的重要表述。}\label{ch16:fig:s30-21-4}
\end{minipage}\par\medskip
图~\ref{ch16:fig:s30-21-4}按卡诺、开尔文、克劳修斯的顺序排列，分别对应前述热机研究及两种早期陈述。

\section{熵与热力学第二定律}

在得到热力学第二定律的早期陈述以后，人们逐渐认识到：对于任意热力学系统，除压强、密度、温度、
内能等状态变量外，还存在另一个状态变量——熵（entropy），记为 $S$。熵从出现之初就带有神秘色彩，
尤其是由熵增出发对宇宙结局的推测；正因为如此，更应准确理解熵的定义和适用条件，避免在不正确的
条件下滥用相关定律。

熵首先是状态变量。与内能一样，它完全由热力学系统的状态决定，而与系统以什么过程到达该状态无关。
状态变量必然同组成系统的微观粒子的性质有关：压强是粒子碰撞容器表面的宏观效果；在经典图像中，
温度由粒子的平均能量决定；内能可理解为全部微观粒子的能量总和。熵也由微观状态决定，玻尔兹曼后来
揭示了这一联系。

当热力学系统经历某个过程时，熵可能改变。对于可逆过程中的任意状态，都可以定义系统温度。若系统在
温度为 $T$ 的平衡态吸收一小份热量 $\delta Q_{\rm rev}$，且热量足够小，使温度在这一微元过程中可视
为不变，则熵的改变为
\begin{equation}
  \boxed{\dd S=\frac{\delta Q_{\rm rev}}{T}}.
\end{equation}
可逆吸热时熵增加，可逆放热时熵减小；相同热量在较低温度下传递，引起的熵变更大。可逆绝热过程没有
热量交换，因而也是等熵过程。与内能相似，经典热力学首先给出的是熵在过程中的改变，而不是任意系统
熵的绝对数值。

不可逆过程同样会导致熵变。由于熵是状态变量，虽然不能把不可逆路径上的每一中间状态都当作平衡态，
仍可另选一条连接同一初、末态的可逆路径来计算 $\Delta S$。对孤立系统，任意自发过程只可能使熵增加；
在理想可逆极限中熵保持不变，孤立系统中的过程不能使总熵减小。这就是熵增原理。其直接推论是：孤立
系统的热平衡态是熵最大的状态。

在熵概念提出的早期，克劳修斯曾把这一结论直接应用于整个宇宙，设想宇宙最终达到处处均匀、温度相同
的热平衡状态，所有细节全部消失。今天我们知道，不能把一个简单孤立热力学系统的结论不加条件地套用
到演化的宇宙；宇宙中丰富结构的形成还受到引力、膨胀和非平衡过程等因素控制。

\subsection{准静态为什么还不够}

准静态只表示变化足够缓慢，使系统能用一串近似平衡态描述；可逆还要求系统和环境都能在不留下
其他变化的条件下复原。例如，缓慢拖动有摩擦的活塞可以近似准静态，但摩擦仍把机械能转为内能，
单纯反向拖动不能把这些内能全部恢复为原来的机械能。有限温差传热也有同样的问题。
因此，可逆传热须令两侧温差趋于零，可逆体积功须令内外压强差趋于零，并排除摩擦等耗散。
这是一种极限模型，不是要求实际装置在有限时间内严格达到它。

\subsection{熵流与熵产生}

对不交换物质的系统，热量可以从不同温度的环境输入。设第 $j$ 个恒温热库温度为 $T_j$，
在全过程向系统传入的带符号热量为 $Q_j$。热库自身的熵变为 $-Q_j/T_j$。
把系统和这些热库一起考察，第二定律要求总熵变非负，因而
\begin{equation}
 \Delta S_{\rm sys}-\sum_j\frac{Q_j}{T_j}=S_{\rm gen}\ge0.
\end{equation}
这里 $S_{\rm gen}$ 是系统连同有关传热边界的总熵产生；它不是新的能量。
等号对应整个过程可逆。如果环境温度变化，应把 $Q_j/T_j$ 换成相应的积分，不能随意选一个平均温度。
对循环，$\Delta S_{\rm sys}=0$，于是
\begin{equation}
 \sum_j\frac{Q_j}{T_j}\le0.
\end{equation}
这就是以恒温热库分段表示的克劳修斯不等式。循环可以把吸热和放热的熵交换相互抵消，
却不能用负的熵产生去抵消耗散。

例如，两热库间直接传递热量 $Q>0$，高温为 $T_h$、低温为 $T_c<T_h$。
高温热库熵变 $-Q/T_h$，低温热库熵变 $Q/T_c$，故
\[
 S_{\rm gen}=Q\left(\frac1{T_c}-\frac1{T_h}\right)>0.
\]
总能量没有损失，但可以利用的温差减小了。若 $T_h\to T_c$，单位热量引起的熵产生趋于零，
这正是可逆传热必须取无穷小温差的原因。

\subsection{由熵不等式判断热机与制冷机}

在 $T_h>T_c>0$ 的两个热库之间，热机每循环从高温热库取热 $Q_h>0$，向低温热库排热 $Q_c>0$。
工质循环复原，环境总熵变为
\[
 -\frac{Q_h}{T_h}+\frac{Q_c}{T_c}\ge0
 \quad\Longrightarrow\quad
 \frac{Q_c}{Q_h}\ge\frac{T_c}{T_h}.
\]
结合 $W=Q_h-Q_c$，得到
\begin{equation}
 \eta=1-\frac{Q_c}{Q_h}\le1-\frac{T_c}{T_h}.
\end{equation}
这次推导没有使用气体物态方程，所以卡诺效率是任意工质的上限，而不只是理想气体计算的巧合。
若装置的某一部分不可逆，通常有严格小于号。

制冷机的有用效果是从低温热库取出 $Q_c$，外界输入功 $W>0$，向高温热库排出
$Q_h=Q_c+W$。此时热库熵变满足
\[
 \frac{Q_h}{T_h}-\frac{Q_c}{T_c}\ge0
 \quad\Longrightarrow\quad
 W\ge Q_c\frac{T_h-T_c}{T_c}.
\]
制冷系数和供热系数分别为
\begin{equation}
 \varepsilon_{\rm ref}=\frac{Q_c}{W}\le\frac{T_c}{T_h-T_c},\qquad
 \varepsilon_{\rm hp}=\frac{Q_h}{W}=1+\varepsilon_{\rm ref}
 \le\frac{T_h}{T_h-T_c}.
\end{equation}
二者可以大于1，因为它们比较的是搬运热量与输入功，而不是把输入热量转为功的比例。
例如 $T_h=300\,\mathrm K,T_c=270\,\mathrm K$ 时，可逆制冷系数为9；搬走
$900\,\mathrm J$ 的热量至少需要 $100\,\mathrm J$ 的功，向高温环境排热为
$1000\,\mathrm J$。由 $1000/300=900/270$ 可同时核验熵平衡。


\section{具有简单物态方程系统的熵}

在绝大多数情形下，热力学系统的熵不容易直接求出；但对物态较简单的系统，某些过程中的熵变可以精确
或近似计算。最简单的系统之一是热库。它有极大的热容量，无论吸收还是放出所研究的有限热量，其温度
都可以近似认为不变。因此，温度为 $T$ 的热库吸收热量 $Q$ 时熵变就是 $Q/T$。

另一类系统以固体和液体为代表。通常吸、放热会引起温度变化，虽然物质体积也随温度改变，但固体、液体
的体积变化往往很小。若忽略体积变化，就同时忽略了系统与环境通过体积功发生的能量交换。设系统比热
为 $c$，并在所研究温区内近似为常数。质量为 $m$ 的物质在温度为 $T$ 时可逆吸收
$\dd Q=mc\,\dd T$，因而
\[
  \dd S=\frac{\dd Q}{T}=mc\frac{\dd T}{T}.
\]
从参考温度 $T_0$ 积分到任意温度 $T$，得到
\begin{equation}
  \boxed{S(T)-S(T_0)=mc\ln\frac{T}{T_0}}.
\end{equation}

理想气体的熵不仅依赖温度，还依赖体积。设有 $n$ 摩尔理想气体，参考状态为 $(V_0,T_0)$，要计算任意
状态 $(V,T)$ 相对参考状态的熵变，先假设 $c_v$ 在温区内为常数。由于熵是状态量，可以任意设计一条连接两状态的可逆路径。图~\ref{ch16:fig:s30-21-5}
选择两步路径：先作等容过程 $(V_0,T_0)\to(V_0,T)$，再作等温过程 $(V_0,T)\to(V,T)$。第一步熵变为
\[
  S(V_0,T)-S(V_0,T_0)=n c_v\ln\frac{T}{T_0};
\]
第二步中理想气体的吸热量等于对外功，故
\[
  S(V,T)-S(V_0,T)=nR\ln\frac{V}{V_0}.
\]
两式相加即得
\begin{equation}
  \boxed{S(V,T)-S(V_0,T_0)
  =n c_v\ln\frac{T}{T_0}+nR\ln\frac{V}{V_0}}.
\end{equation}

\par\noindent\begin{minipage}{\linewidth}
\entropyPathFigure
\diagramcaption{理想气体任意两状态间的熵差。图中用等容升温加等温膨胀的可逆路径代替任意真实路径。}\label{ch16:fig:s30-21-5}
\end{minipage}\par\medskip
图~\ref{ch16:fig:s30-21-5}的 $A\to B$ 改变温度而不改体积，$B\to C$ 改变体积而不改温度；因此两个对数项各有清楚的来源。
若实际过程是降温或压缩，相应积分上限小于下限，对数自动给出负号，并不要求实际过程也按图运行。
若 $c_v$ 随温度变化，应保留积分
\[
 \Delta S=n\int_{T_0}^{T}\frac{c_v(T')}{T'}\,\dd T'+nR\ln\frac{V}{V_0}.
\]

\subsection{熵变用初末态计算，熵产生还要看实际环境}
选一条可逆辅助路径，是为了计算状态函数的差，不是把真实过程宣称为可逆。图\ref{ch16:fig:s30-21-5}中的 $A\to B\to C$ 正是这种计算工具。系统熵变、热库熵变和二者的总熵变须分开；系统自己可以减熵，只要包括有关环境后的总熵不减。

\Needspace{180pt}
\begin{origexamplebox}{同末态的两种膨胀，熵产生为何不同}\label{ch16:ex:teaching-two-paths}
$n$ 摩尔理想气体初态为 $(T_0,V_0)$。过程甲是在 $T_0$ 热库作用下可逆等温膨胀至 $2V_0$；过程乙是在隔热刚性容器中向真空自由膨胀至 $2V_0$。求气体、有关环境的熵变和总熵产生。忽略容器热容量、整体动能及势能的变化。

\ThQTwoPaths
\end{origexamplebox}
\Needspace{65pt}
\begin{xesolution}
\ThATwoPaths

乙过程既不传热也不向外做功，内能不变，理想气体末温仍为 $T_0$。所以甲、乙的气体初末态完全相同，均有
\[
 \Delta S_{\rm gas}=nc_v\ln\frac{T_0}{T_0}+nR\ln2=nR\ln2.
\]
甲从热库吸热 $Q=nRT_0\ln2$，热库熵变 $-Q/T_0=-nR\ln2$。接收功的理想机械装置不增加熵，故甲的总熵产生为零。
乙没有向环境传热，环境熵变为零，因此
\[
 S_{{\rm gen},\text{甲}}=0,\qquad
 S_{{\rm gen},\text{乙}}=nR\ln2>0.
\]
乙中 $Q=0$ 却有 $\Delta S_{\rm gas}>0$，并不违反熵定义，因为 $\Delta S=\int\delta Q_{\rm rev}/T$ 中的热量属于辅助可逆路径，而非真实自由膨胀的热量。若末体积趋于初体积，两者熵变都趋零；$nR$ 的单位为 $\mathrm{J/K}$，对数中的体积比无量纲。
\end{xesolution}

\subsection{有限热源不能始终当作恒温热库}
实际物体一边供热一边降温，热机每一小步可达到的效率也在改变。因此不能选一个平均温度代入卡诺效率再乘总热量。解决有限热源问题的顺序是：先明确装置是否循环复原，再写总能量和总熵条件，最后判断上界能否由可逆过程逼近。

\Needspace{180pt}
\begin{origexamplebox}{两物体之间能够取出的最大功}\label{ch16:ex:teaching-finite-work}
两个物体的恒定热容量为 $C_1,C_2$，初温为 $T_h>T_c>0$。在它们之间运行热机，直到两者达到共同终温，热机最后恢复原态。整体与外界不交换热量，但可以输出机械功；忽略体积变化和其他能量改变。求最多能输出多少功，并与两物体直接接触比较。

\ThQFiniteWork
\end{origexamplebox}
\Needspace{65pt}
\begin{xesolution}
\ThAFiniteWork

若共同终温为 $T_f$，热机不储存净能量，输出功为
\[
 W=C_1T_h+C_2T_c-(C_1+C_2)T_f.
\]
能量零点抵消，只涉及温度差。总熵不能减少：
\[
 C_1\ln\frac{T_f}{T_h}+C_2\ln\frac{T_f}{T_c}\ge0.
\]
把第一项的对数拆成 $\ln(T_f/T_c)-\ln(T_h/T_c)$，得到
\[
 (C_1+C_2)\ln\frac{T_f}{T_c}\ge C_1\ln\frac{T_h}{T_c},\qquad
 T_f\ge T_*=T_c\left(\frac{T_h}{T_c}\right)^{C_1/(C_1+C_2)}.
\]
所有对数和幂函数的自变量均无量纲。由于 $W$ 随 $T_f$ 增加而减小，最大功为
\[
 W_{\max}=C_1(T_h-T_*)+C_2(T_c-T_*).
\]
可在两物体的瞬时温度之间运行无穷小可逆循环，使总熵保持不变，直到温差趋零，所以该上界可以在可逆极限下逼近。图\ref{ch16:fig:s30-21-3}中的普通热机能流方向仍适用，但两端温度在不断变化。
直接接触时不输出功，终温是加权算术平均 $T_d=(C_1T_h+C_2T_c)/(C_1+C_2)\ge T_*$。若 $C_1=C_2=C$，则
\[
 T_*=\sqrt{T_hT_c},\qquad W_{\max}=C(\sqrt{T_h}-\sqrt{T_c})^2.
\]
两初温相等时最大功为零。直接传热没有损失总能量，却损失了温差所能提供的做功能力；不能把“能量守恒”误读为“有用功不变”。
\end{xesolution}

\Needspace{180pt}
\begin{origexamplebox}{将有限物体冷却到环境温度以下所需的最少功}\label{ch16:ex:teaching-cooling-work}
热容量为常数 $C$ 的物体初温与环境热库同为 $T_0$。用循环复原的制冷机将物体冷却到 $T_1<T_0$，环境保持 $T_0$。求最少输入功。忽略体积变化及漏热，假定此温区无相变。

\ThQCooling
\end{origexamplebox}
\Needspace{65pt}
\begin{xesolution}
\ThACooling

在物体瞬时温度 $T$ 下再降低 $\dd T<0$，须取出热量 $\delta Q_c=-C\dd T>0$。可逆制冷系数为 $T/(T_0-T)$，所以最小元功为
\[
 \delta W_{\min}=(-C\dd T)\frac{T_0-T}{T}.
\]
系数随温度改变，必须积分：
\[
 W_{\min}=C\int_{T_1}^{T_0}\left(\frac{T_0}{T}-1\right)\dd T
 =C\left[T_0\ln\frac{T_0}{T_1}-(T_0-T_1)\right].
\]
物体总计放出 $Q_c=C(T_0-T_1)$，环境得到 $Q_h=Q_c+W_{\min}=CT_0\ln(T_0/T_1)$。因此物体熵变 $C\ln(T_1/T_0)$ 恰被环境的 $Q_h/T_0$ 抵消，验证了可逆极限。
令 $x=(T_0-T_1)/T_0$，小温差时 $W_{\min}\simeq CT_0x^2/2>0$，而不是与温差一次方成正比。真实有限速率制冷还会产生熵，所需功更大；本模型不能跨越热容变化或相变而直接外推至绝对零度。
\end{xesolution}

\section{熵的微观解释}

宏观量内能可以解释为微观粒子动能与势能的总和。熵也最初来自宏观热过程的研究，但随着分子运动理论
的发展，人们逐渐认识到它的微观含义。玻尔兹曼最早系统提出基于分子运动理论的熵解释。虽然这一理论
起初饱受争议，但后来被广泛接受；在主流学界尚未普遍接受玻尔兹曼理论时，普朗克和爱因斯坦已经
基于其统计思想建立早期量子理论。

\par\noindent\begin{minipage}{\linewidth}
\boltzmannMemorialFigure
\diagramcaption{玻尔兹曼纪念碑上的熵公式，把宏观熵与微观状态的可能数目联系起来。}\label{ch16:fig:s30-21-6}
\end{minipage}\par\medskip
图~\ref{ch16:fig:s30-21-6}中纪念碑上方的公式，是接下来计数解释的核心。

今天通常把玻尔兹曼公式写成
\begin{equation}
  \boxed{S=k_B\ln\Omega},
\end{equation}
其中 $k_B$ 是玻尔兹曼常数，$\Omega$ 是给定宏观状态所对应的微观状态数。一个宏观状态允许的微观
排布可能性越多，熵就越大；这与把熵理解为系统“无序程度”的传统图像接近，但微观状态数的说法更准确。

先看气体分布这一例子。把容器分为等大的左、右两半，暂只记录每个粒子位于哪一半，而不记录具体位置和速度。
给 $N$ 个粒子临时编号，是为了计数各粒子的左右选择，不是把编号当作新的物理自由度。
若左侧有 $N_1$ 个，右侧有 $N-N_1$ 个，则先从全部粒子中挑出左侧粒子的有序序列，有
$N(N-1)\cdots(N-N_1+1)$ 种；同一组左侧粒子的内部排列被重复计了 $N_1!$ 次，故应除去：
\[
 \Omega_{\rm LR}(N_1)=\frac{N(N-1)\cdots(N-N_1+1)}{N_1!}
 =\frac{N!}{N_1!(N-N_1)!}.
\]
例如 $N=4$ 时，全在左侧仅1种，两边各2个有6种。一般有
$\Omega_{\rm LR}(N_1+1)/\Omega_{\rm LR}(N_1)=(N-N_1)/(N_1+1)$，
所以计数先增后减，在左右数目最接近时最大。每个粒子左右等可能且彼此独立时，
概率为 $\Omega_{\rm LR}/2^N$。这解释了为什么宏观气体倾向于分散到整个容器。
$\Omega_{\rm LR}$ 只计算这一粗粒化的分布因素，并不是包含全部位置、动量的完整微观状态数。

同一组合计数也可用于顺磁性物质。某些原子具有非零磁矩，其性质有点像指南针；在外磁场作用下，简化模型中每个磁矩
只有沿磁场或逆磁场两个方向。设系统含 $N$ 个原子，其中 $N_1$ 个磁矩向上，剩下
$N_2=N-N_1$ 个向下。宏观测量只能确认上下数目的差，不能确定究竟哪一些原子向上。于是，在 $N$ 个
原子中选出 $N_1$ 个向上的微观排布数为
\begin{equation}
  \boxed{\Omega=\frac{N!}{N_1!(N-N_1)!}}.
\end{equation}
在没有外场时，$N_1=N/2$ 附近的微观状态数最大，这正是通常情况下顺磁性物质不显示宏观磁性的统计
原因。

中学阶段通常不会进行复杂的微观熵运算，但必须注意微观世界与宏观世界之间巨大的数目差异。设一个人
吃下一份标称化学能为 $\SI{40}{J}$ 的食物，体温为 $\SI{310}{K}$，能否把人体熵增算成 $40/310$？
不能：进食是物质流入和化学反应，不等于向恒温封闭系统传入40焦耳的热量；还要考虑物质携熵、散热和代谢。
为比较数目大小，改考察一个明确的理想模型：温度恒为 $\SI{310}{K}$ 的热库净吸收
$\SI{40}{J}$ 热量。这个热库的熵变才是
\[
  \Delta S=\frac{Q}{T}
  =\frac{\SI{40}{J}}{\SI{310}{K}}
  \approx\SI{0.13}{J.K^{-1}}.
\]
另一方面，玻尔兹曼公式给出
\[
  \Delta S=k_B\ln\frac{\Omega_f}{\Omega_i}.
\]
联立两式，取 $k_B=\SI{1.380649e-23}{J.K^{-1}}$，得到
\begin{equation}
  \boxed{\ln\frac{\Omega_f}{\Omega_i}
  =\frac{40/310}{1.380649\times10^{-23}}
  \approx9.35\times10^{21}}.
\end{equation}
这里先报告对数，避免把数量级粗略相近的巨大指数误写成状态数之比的数值近似。按这一统计描述，
该比值可写为 $\exp(9.35\times10^{21})$ 的量级说明；宏观上很小的热量也可对应极大的微观可能数目变化。


\Needspace{12\baselineskip}\section{综合例题}
\begin{origexamplebox}{【卡诺热机】}\label{ch16:ex:s30-21-1}
利用熵增原理证明卡诺定理：任意工作物质作卡诺循环的效率都只与高、低温热库的温度有关，
与工质的数量、物态方程均无关，效率均为
\[
  \eta=1-\frac{T_2}{T_1},
\]
其中 $T_1,T_2$ 分别为高、低温热库的温度。

\ThQCarnot
\end{origexamplebox}\par\nopagebreak[4]

\par\nopagebreak[4]\begin{xesolution}
设热机在一个可逆循环中从温度为 $T_1$ 的高温热库吸收热量 $Q_1$，向温度为 $T_2$ 的低温热库放出
热量 $Q_2$。工质经历完整循环后回到初态，作为状态函数的熵没有净改变。

高温热库失去热量，熵变为 $-Q_1/T_1$；低温热库得到热量，熵变为 $Q_2/T_2$。卡诺循环是可逆
循环，把工质、两热库和只储存机械能的理想功库一并作为孤立整体；功库熵不变，所以总熵变为零，因此
\[
  -\frac{Q_1}{T_1}+\frac{Q_2}{T_2}=0,
  \qquad
  \frac{Q_2}{Q_1}=\frac{T_2}{T_1}.
\]

\par\noindent\begin{minipage}{\linewidth}
\carnotProofDiagram
\figcap{卡诺热机的能量流与熵平衡。$Q_1,Q_2,W$ 均表示图中箭头方向上的正量。}\label{ch16:fig:carnot-proof}
\end{minipage}\par\medskip
图\ref{ch16:fig:carnot-proof}中热量的大小不同，但可逆情况下从高温热库取走的熵等于送给低温热库的熵。

热机一个循环的对外净功为 $W=Q_1-Q_2$，故
\[
  \boxed{\eta=\frac{W}{Q_1}
  =1-\frac{Q_2}{Q_1}
  =1-\frac{T_2}{T_1}}.
\]
推导中没有使用工质的物态方程、热容量或物质的量，只使用了“循环后工质熵不变”和“整个过程可逆”。
所以任意工质作卡诺循环都具有同一效率。若过程不可逆，总熵会增加，向低温热库排出的热量更多，
实际效率必低于卡诺效率。
\end{xesolution}


\Needspace{12\baselineskip}\begin{origexamplebox}{【熵增原理】}\label{ch16:ex:s30-21-2}
利用熵增原理，重新对热力学第二定律的开尔文和克劳修斯陈述给出定量说明。

\ThQSecondLaw
\end{origexamplebox}\par\nopagebreak[4]

\par\nopagebreak[4]\begin{xesolution}
\textbf{开尔文表述。} 假设一台循环装置从温度为 $T$ 的单一热库吸收热量 $Q$，并把它全部变为功。
装置完成循环后熵变为零；热库失去热量，熵变为
\[
  \Delta S_{\rm reservoir}=-\frac{Q}{T}<0.
\]
还需把接收机械功的理想功库纳入整体，例如无摩擦地被提高的重物；它只增加机械能，熵不变。
于是孤立的“装置、热库加功库”总熵减少，违背熵增原理。因此不可能在没有其他变化的条件下，从单一热库
吸热并把热量全部转化为功。

\textbf{克劳修斯表述。} 假设一台循环装置不消耗功，却把热量 $Q$ 从温度 $T_c$ 的低温热库搬到温度
$T_h$ 的高温热库，其中 $T_h>T_c$。两热库的总熵变为
\[
  \Delta S_{\rm tot}
  =\frac{Q}{T_h}-\frac{Q}{T_c}<0.
\]
这同样违背熵增原理。因此把热量从低温处送往高温处，必须伴随外界做功或其他补偿变化。

\par\noindent\begin{minipage}{\linewidth}
\secondLawEntropyDiagram
\figcap{两种被禁止的唯一效果及其负熵变。比较时都须把装置和所有环境变化计入。}\label{ch16:fig:second-law-comparison}
\end{minipage}\par\medskip
图\ref{ch16:fig:second-law-comparison}左侧的热机必须已完成循环；单次等温膨胀并不属于这个禁止条件。

两个陈述表面上分别谈“热变功”和“逆向传热”，本质上都排除了使孤立系统总熵减少的循环装置。
\end{xesolution}


\Needspace{12\baselineskip}\begin{origexamplebox}{【统统证出来】}\label{ch16:ex:s30-21-3}
请利用热力学第二定律辨析下列说法，指出所需条件，并给出正确结论：
\begin{enumerate}[label=\alph*.]
  \item 理想气体的等温膨胀过程是不可能自然发生的；
  \item 对本讲的定量理想气体，可逆绝热线与等温线不能相交两次；
  \item 对固定组成、稳定单相的简单可压缩系统，两条不同可逆绝热线不能相交；
  \item 有限温差热传递过程不可逆；
  \item 快速膨胀过程不可逆。
\end{enumerate}

\ThQIrreversibility
\end{origexamplebox}\par\nopagebreak[4]

\par\nopagebreak[4]\begin{xesolution}
这组判断的共同方法是：若假定某过程可以在没有其他影响的情况下逆转，就尝试拼出第二类永动机，或
直接算出总熵减少。

\textbf{a. 等温膨胀。} 这一说法本身不成立，不能给“自然发生”附加原本没有的限制来证明它。
理想气体在恒温环境中推动活塞缓慢膨胀可以发生，单个过程并不违反第二定律。此时 $\Delta U=0$，对外功
\[
  W=nRT\ln\frac{V_2}{V_1}>0
\]
必须完全来自热库供热。工质的体积已经改变，这不是“装置循环复原后只留下热变功”这一被禁止的效果。
它不能在既绝热又输出功的条件下仍维持等温，这首先可由第一定律判断。另一个更直接的反例是撤去
隔板，让理想气体向真空自由膨胀，则 $Q=W=\Delta U=0$，温度确实不变，而且过程会自发发生；这正是
不可逆等温自由膨胀，不能与可逆等温膨胀混为一谈。

\par\noindent\begin{minipage}{\linewidth}
\irreversibleProcessesDiagram
\figcap{左侧分别画出撤去隔板前、后的气体状态；右侧为两恒温热库间的有限温差传热。}\label{ch16:fig:irreversibility}
\end{minipage}\par\medskip
图\ref{ch16:fig:irreversibility}中左图初态有真空区，末态虚线只标隔板原位置，不再是阻挡气体的实体壁；
右图的传热将在第d问中计算总熵变。

\textbf{b. 绝热线与等温线的交点。} 对本讲的理想气体，固定温度下
$S(T,V)=\text{常量}+nR\ln V$ 随 $V$ 严格增加，而同一可逆绝热线有固定熵，故只能选出一个 $V$。
也可代入 $P=nRT/V$ 与 $PV^\gamma=K$，得 $nRTV^{\gamma-1}=K$，在 $\gamma>1,V>0$ 时解唯一。
若强行假设两次相交并能围成非零功的循环，就可用一段等温线与一段绝热线构造单热库循环，违反第二定律。
对于一般稳定单相系统，还需可逆绝热线上的温度随体积单调，例如通常热膨胀系数为正的介质。
仅有正热容不能保证这个单调性；存在反常热膨胀时，不能把“两次相交不可能”无条件推广。

\textbf{c. 两条绝热线不能相交。} 可逆绝热线就是等熵线。若两条绝热线在某平衡状态相交，它们在该点
具有相同的熵；对于给定组成、单相稳定且物态方程光滑的简单系统，$c_v>0$ 使
$(\partial S/\partial T)_V=nc_v/T>0$，故固定熵和体积可唯一确定温度。通过该点的可逆绝热路径由微分方程唯一确定，因此
两条曲线实际是同一条绝热线，而不是两条不同曲线。若把它们看作不同曲线，还可以在两曲线间补上一小段
等温线，构造出上一问所排除的单热库循环。

\textbf{d. 有限温差传热不可逆。} 先设两个恒温热库，热量 $Q$ 从 $T_h$ 传到 $T_c$，$T_h>T_c$。总熵变
\[
  \Delta S_{\rm tot}
  =-\frac{Q}{T_h}+\frac{Q}{T_c}
  =Q\left(\frac1{T_c}-\frac1{T_h}\right)>0.
\]
若要完全逆转，就必须使孤立系统熵减少，故不能自发实现。只有在传热温差趋于零的可逆极限中，熵产生
才趋于零。
若两个有限物体的温度会变化，则应对微小传热量积分：
$\Delta S_{\rm tot}=\int(1/T_c-1/T_h)\,\delta Q>0$；不能把初始温度直接代入有限热量公式。

\textbf{e. 快速膨胀。} “快”本身不是可直接代入熵公式的物理条件。通常快速膨胀来不及维持平衡，
伴随有限压差、黏滞、冲击或最终弛豫时，会产生熵，因而不可逆；若只给一个短时间而没有机制，证明并不充分。
以绝热容器中撤去隔板、气体向真空自由膨胀为明确模型：气体由 $V_i$ 扩展到 $V_f$，温度不变，但
\[
  \Delta S=nR\ln\frac{V_f}{V_i}>0.
\]
若想让气体与环境都无变化地回到初态，就必须使总熵减少，所以过程不可逆。

\end{xesolution}


\Needspace{12\baselineskip}\begin{origexamplebox}{【热传递】}\label{ch16:ex:s30-21-4}
三个热容量均为 $c$ 的物体，已知 $T_1=T_2=\SI{300}{K}$，$T_3=\SI{100}{K}$。若外界不做功也不供热，
利用热机工作于三物体之间，使三物中某一物体温度升高，求它所能达到的最高温度；此时另外两物体的
温度各为多少？

\ThQThreeBodies
\end{origexamplebox}\par\nopagebreak[4]

\par\nopagebreak[4]\begin{xesolution}
先证明任何允许过程的温度上界，再证明这个上界可由可逆装置趋近，不能只用两条等式猜测最优状态。
取最终最高温度为 $X$，另外两体温度为 $Y,Z$，均以开尔文计。热机完成循环，不储存净能量，
则第一定律和熵增原理分别给出
\[
 X+Y+Z=700,\qquad
 \Delta S=c\ln\frac{XYZ}{300^2\times100}\ge0.
\]
所以 $XYZ\ge9.0\times10^6$。又由 $(Y-Z)^2\ge0$，
\[
 YZ\le\left(\frac{Y+Z}{2}\right)^2=\left(\frac{700-X}{2}\right)^2,
 \quad X(700-X)^2\ge36.0\times10^6.
\]
因为 $X$ 是最高温度，$X\ge700/3$。函数 $f(X)=X(700-X)^2$ 在该区间内有
$f'(X)=(700-X)(700-3X)\le0$；而 $f(400)=36.0\times10^6$，故必有 $X\le400$。
等号同时要求总熵不变及 $Y=Z$。这才说明可逆性和另外两体等温是上界成立的必要条件。

设目标物体的最高温度为 $T_{\max}$。三个物体的热容量均为常量 $c$，系统与外界既不交换热量也不交换
功，因此总能量守恒：
\[
  c(T_{\max}+2T)=c(300+300+100),
  \qquad T_{\max}+2T=700.
\]
可逆极限下总熵不变：
\[
  c\ln\frac{T_{\max}}{300}
  +c\ln\frac{T}{300}
  +c\ln\frac{T}{100}=0.
\]
于是
\[
  T_{\max}T^2=300\times300\times100.
\]

\par\noindent\begin{minipage}{\linewidth}
\threeBodiesDiagram
\figcap{两台可逆装置在三物体之间协同工作：热机的功全部驱动热泵，不从外界取得净功。}\label{ch16:fig:three-body-engine}
\end{minipage}\par\medskip

联立 $T_{\max}=700-2T$，得到
\[
  (700-2T)T^2=9.0\times10^6.
\]
物理解为 $T=\SI{150}{K}$，从而
\[
  \boxed{T_{\max}=\SI{400}{K}},
  \qquad
  \boxed{T_{\rm other}=\SI{150}{K},\ \SI{150}{K}}.
\]
方程还有 $T=\SI{300}{K},T_{\max}=\SI{100}{K}$ 这一根，它只是把初态中三物体的编号重新排列，并非
所求的最高温度状态。
最后验证可实现性。图\ref{ch16:fig:three-body-engine}中，令供热体 $Y$ 与低温体 $Z$ 之间的可逆热机输出功，
恰好驱动从 $Z$ 向目标体 $X$ 输热的可逆热泵。每小步都令两机完成循环，故没有外界净功输入。
热机从 $Y$ 取出 $-c\dd Y>0$，输出功
$\delta W=(-c\dd Y)(1-Z/Y)$；热泵向 $X$ 送入 $c\dd X$，需要功
$\delta W=c\dd X(1-Z/X)$。二者相等，等价于
\[
 \frac{\dd X}{X}+\frac{\dd Y}{Y}+\frac{\dd Z}{Z}=0,
 \qquad \dd X+\dd Y+\dd Z=0.
\]
沿这一可逆路径，$XYZ=9.0\times10^6$ 且 $X+Y+Z=700$ 始终成立。对 $300\le X<400$，
另外两个正根是
\[
 Y,Z=\frac{700-X\pm\sqrt{(700-X)^2-36.0\times10^6/X}}{2},
\]
并满足 $X\ge Y>Z$，热机与热泵均可按图运行。随着 $X\to400$，$Y,Z\to150$，
两者的温差趋于零，达到可逆极限；实际有限速率装置只能逼近这一最高温度。
所有对数都作用于温度比，温度积的两边单位都是 $\mathrm K^3$。
\end{xesolution}


\Needspace{12\baselineskip}\begin{origexamplebox}{【液体表面能】}\label{ch16:ex:s30-21-5}
小明学习了卡诺热机后，想把液体表面作为工作物质。液体表面能够对外界物体施加表面张力，因此
表面具有能量，并且可以得到
\[
  U=\left(\sigma-T\frac{\dd\sigma}{\dd T}\right)A,
\]
其中 $\sigma$ 为液体的表面张力系数，$A$ 为液体表面积。请研究利用液体表面作为工作物质的热机，并
求出其效率。

\ThQSurfaceEngine
\end{origexamplebox}\par\nopagebreak[4]

\par\nopagebreak[4]\begin{xesolution}
先约定 $A$ 是所有被计入的气液界面的总面积；若是肥皂膜，须把两面都计入，不能再额外乘2。
本题 $U$ 指随界面增加而产生的表面过剩内能，不包括整块液体的体相内能。
可逆扩展表面时外界对表面做元功 $\delta W^{\rm in}=\sigma\dd A$，因而
\[
 \dd U=T\dd S+\sigma\dd A,\qquad
 \delta W^{\rm out}=-\sigma\dd A.
\]
定义亥姆霍兹（Helmholtz）自由能 $F=U-TS$，有 $\dd F=-S\dd T+\sigma\dd A$。
在忽略曲率效应、$\sigma$ 仅随温度变化时，表面过剩自由能为
\[
  F=\sigma(T)A.
\]
又因为 $F=U-TS$，把题给的 $U$ 代入可得表面熵
\[
  S=\frac{U-F}{T}
  =-A\frac{\dd\sigma}{\dd T}.
\]
普通液体的表面张力通常随温度升高而减小，故 $\dd\sigma/\dd T<0$；在给定温度下扩大表面会增加熵，
并需要从热库吸热。

让液体表面工作在温度分别为 $T_h,T_c$ 的两个热库之间，组成两段等温过程和两段等熵过程。在高温
$T_h$ 下，表面积从 $A_1$ 可逆增加到 $A_2$，表面吸热
\[
  Q_h=T_h\Delta S,
  \qquad
  \Delta S=-\left.\frac{\dd\sigma}{\dd T}\right|_{T_h}(A_2-A_1)>0.
\]
然后沿等熵过程降到 $T_c$，在低温下使表面积减小并向低温热库放热。由于两条等温段的熵变量大小相等，
放出的热量为
\[
  Q_c=T_c\Delta S.
\]
所以一个循环的净功与效率为
\[
  W_{\rm out}=Q_h-Q_c=(T_h-T_c)\Delta S,
  \qquad
  \boxed{\eta=1-\frac{T_c}{T_h}}.
\]

\par\noindent\begin{minipage}{\linewidth}
\surfaceEngineDiagram
\figcap{表面工质的可逆循环。上下为等温段，两条弯曲边为等熵段；面积与温度都随等熵过程改变。}\label{ch16:fig:surface-cycle}
\end{minipage}\par\medskip

还需确定等熵线，不能直接把它画成等面积直线。对 $S=-A\sigma'(T)$ 微分得
\[
 \dd S=-\sigma'(T)\dd A-A\sigma''(T)\dd T=0,
 \qquad A\sigma'(T)=\text{常量}.
\]
一般等熵过程既改变面积也改变温度。图\ref{ch16:fig:surface-cycle}采用一个明确的非退化模型
$\sigma=\sigma_*-aT^2$，其中 $a>0$，温区内 $\sigma>0$。这时
$S=2aAT$、$U=(\sigma_*+aT^2)A$、$C_A=2aAT>0$，等熵线为 $AT=\text{常量}$。
高温等温段 $1\to2$ 的面积为 $A_1\to A_2$；冷却到低温的状态3满足
$A_3=A_2T_h/T_c$；低温压缩到状态4时 $A_4=A_1T_h/T_c$，再沿等熵线回到状态1。
由此 $Q_h=2aT_h^2(A_2-A_1)$、$Q_c=2aT_hT_c(A_2-A_1)$，再次给出相同效率。

如果只在很窄温区用线性表面张力近似，形式上 $\sigma'$ 为常量，等熵线会变成 $A=\text{常量}$，
循环呈矩形，形式功为
\[
  W_{\rm out}=(\sigma_c-\sigma_h)(A_2-A_1).
\]
但仅含表面项时这还导致 $C_A=-TA\sigma''=0$，是退化近似，不能据此描述真实液膜如何绝热升降温。
若计入体相自由能 $F_0(T)$，总熵变成 $S_0(T)-A\sigma'(T)$，等熵线必须重新求，通常不再竖直。
无论哪一可靠模型，只要可逆且工作在两个热库之间，效率仍为 $1-T_c/T_h$。
在 $\sigma$--$A$ 图中，本题是逆时针循环，因 $W_{\rm out}=-\oint\sigma\dd A$，其净功为正；
这与气体的 $W_{\rm out}=\oint P\dd V$ 符号相反，不能只凭“循环方向”判断正负。
\end{xesolution}
}

\endgroup

\hypersetup{bookmarksdepth=2}
